\section*{ID8226: Executive Summary (One Page): S\textasciicircum{}(X)\_ij,rm EM}
\paragraph{Metadata.}
PiID: 3363948045759 | RTEID: RTE:P31943.3848:T2.3674 | UUID: 65d119ea-1ccd-56be-baa9-ad739f2e64d3

\paragraph{Canonical Statement.}
**Key formulas (ready to use):** - Covariant EM stress-energy \textbackslash{}[ T\textasciicircum{}\{\textbackslash{}mu\textbackslash{}nu\}\_\{\textbackslash{}rm EM\} = F\textasciicircum{}\{\textbackslash{}mu\textbackslash{}alpha\}F\textasciicircum{}\textbackslash{}nu\{\}\_\textbackslash{}alpha - \textbackslash{}tfrac14 g\textasciicircum{}\{\textbackslash{}mu\textbackslash{}nu\} F\textasciicircum{}\{\textbackslash{}alpha\textbackslash{}beta\}F\_\{\textbackslash{}alpha\textbackslash{}beta\}. \textbackslash{}] - Measured energy density and momentum (observer \textbackslash{}(X\textbackslash{}) = \textbackslash{}(n\textbackslash{}) or \textbackslash{}(u\textbackslash{})): \textbackslash{}[ \textbackslash{}mathcal E\textasciicircum{}\{(X)\} = X\_\textbackslash{}mu X\_\textbackslash{}nu T\textasciicircum{}\{\textbackslash{}mu\textbackslash{}nu\},\textbackslash{}qquad S\textasciicircum{}\{(X)\}\_i = -\textbackslash{}Pi\textasciicircum{}\{(X)\}\_\{i\textbackslash{}mu\}X\_\textbackslash{}nu T\textasciicircum{}\{\textbackslash{}mu\textbackslash{}nu\}, \textbackslash{}] with projector \textbackslash{}(\textbackslash{}Pi\textasciicircum{}\{(n)\}=\textbackslash{}gamma,\textbackslash{} \textbackslash{}Pi\textasciicircum{}\{(u)\}=h\textbackslash{}). - Maxwell stress (spatial) \textbackslash{}[ S\textasciicircum{}\{(X)\}\_\{ij,\textbackslash{}rm EM\} = E\textasciicircum{}\{(X)\}\_iE\textasciicircum{}\{(X)\}\_j + B\textasciicircum{}\{(X)\}\_iB\textasciicircum{}\{(X)\}\_j - \textbackslash{}tfrac12 \textbackslash{}Pi\textasciicircum{}\{(X)\}\_\{ij\}\textbackslash{}big(E\textasciicircum{}\{(X)\}\{\}\textasciicircum{}2+B\textasciicircum{}\{(X)\}\{\}\textasciicircum{}2\textbackslash{}big). \textbackslash{}]

\paragraph{Canonical Equation.}
\[
S^{(X)}_{ij,\rm EM} = E^{(X)}_iE^{(X)}_j + B^{(X)}_iB^{(X)}_j - \tfrac12 \Pi^{(X)}_{ij}\big(E^{(X)}{}^2+B^{(X)}{}^2\big).
\]

\paragraph{Dictionary.}
Key formulas (ready to use): - Covariant EM stress-energy - Measured energy density and momentum (observer \textbackslash{}(X\textbackslash{}) = \textbackslash{}(n\textbackslash{}) or \textbackslash{}(u\textbackslash{})): with projector \textbackslash{}( \textasciicircum{}(n)= ,\textbackslash{} \textasciicircum{}(u)=h\textbackslash{}) (S\textasciicircum{}(X)\_ij,rm EM)

\paragraph{Encyclopedia.}
Executive Summary (One Page): S\textasciicircum{}(X)\_ij,rm EM is an algebraic definition, written S\textasciicircum{}(X)\_ij,rm EM = E\textasciicircum{}(X)\_iE\textasciicircum{}(X)\_j + B\textasciicircum{}(X)\_iB\textasciicircum{}(X)\_j - tfrac12 Π\textasciicircum{}(X)\_ij (E\textasciicircum{}(X)\textasciicircum{}2+B\textasciicircum{}(. It is an algebraic definition expressing the quantity directly in terms of the variables on its right-hand side. It is built from Π, E, X, j, B, defining S\textasciicircum{}(X)\_ij,rm EM. Within the Trawin Zero Point registry it serves as one element of the framework's mathematical narrative.
